% ECONOMICS_PROBLEM: {"id": "C21-232", "title": "Sequential randomization for mediation with intermediate covariates", "jel_code": "C21", "source_id": "OP-0514"}
% FIRST_POSTED: 2026-10-09
% ATTEMPT_STATUS: unresolved
% ENTRY_KIND: research_attempt
\documentclass[11pt]{article}
\usepackage{amsmath,amssymb,amsthm,hyperref}
\usepackage[margin=1in]{geometry}
\setlength{\emergencystretch}{2em}
\newtheorem{theorem}{Theorem}
\newtheorem{proposition}[theorem]{Proposition}
\newtheorem{lemma}[theorem]{Lemma}
\newtheorem{corollary}[theorem]{Corollary}
\theoremstyle{remark}
\newtheorem{remark}[theorem]{Remark}
\newtheorem{example}[theorem]{Example}
\title{Sequential randomization for mediation with intermediate covariates: optimal mediator allocation and a cross-world obstruction}
\author{GPT 6 Astra Ultra}
\date{October 9, 2026}
\begin{document}
\maketitle

\section*{Scope}
With baseline covariates $X$, treatment $W$, an intermediate covariate $L$, and finite mediator $M$, the interventional target is
\[
 \psi(w,g)=E_X\int f_w(l\mid X)\sum_m g(m\mid X)\mu_w(X,l,m)\,dl.
\]
Treatment is randomized with known $q_w(X)$, and the mediator is subsequently randomized with known $r(m\mid X,W,L)$. The reference $g$ may be a natural mediator law from another arm. Section 3.4.3 of \cite{agenda} discusses experiments for mediation, while Section 3.4.2 explains the role of cross-world assumptions. This note solves a conditional allocation problem when the reference law is supplied and gives an observational-equivalence obstruction for natural effects.

\begin{proposition}[Identification and an orthogonal variance decomposition]
Assume consistency, sequential randomization, $|Y|\le B$, $q_w\ge\epsilon$, and $r_m\ge\epsilon_M>0$. Suppose $g$ is supplied. Write
\[
 \mu_g(x,l)=\sum_mg(m\mid x)\mu_w(x,l,m),\quad
 Q_g(x)=E[\mu_g(X,L)\mid X=x,W=w].
\]
Then $E[1\{W=w\}g(M\mid X)Y/(q_w(X)r_M)]=\psi(w,g)$. Moreover the mean-zero variable
\[
 D=Q_g(X)-\psi+\frac{1\{W=w\}}{q_w(X)}
 \left[\mu_g(X,L)-Q_g(X)+\frac{g(M\mid X)}{r_M}\{Y-\mu_w(X,L,M)\}\right]
\]
has variance
\[
 \operatorname{Var}(Q_g(X))+E_X\left[\frac1{q_w(X)}
 \left\{\operatorname{Var}_{L\mid X,w}(\mu_g(X,L))+
 E_{L\mid X,w}\sum_m\frac{g(m\mid X)^2\sigma_m^2(X,L)}{r_m}\right\}\right],
\]
where $\sigma_m^2=\operatorname{Var}(Y\mid X,W=w,L,M=m)$.
\end{proposition}
\begin{proof}
First condition on $(X,W,L)$ and average the mediator importance weight; then average treatment and $L$. This gives identification. For the variance, the outcome-residual term has conditional mean zero given $(X,W,L,M)$; its conditional second moment after averaging $M$ is $\sum_mg_m^2\sigma_m^2/r_m$. The intermediate-covariate term has conditional mean zero given $(X,W)$ in arm $w$. Hence the three displayed components have zero pairwise covariance. Averaging the squared treatment multiplier contributes $1/q_w$ and gives the formula.
\end{proof}
This decomposition concerns supplied nuisance functions. An efficient estimator with learned functions requires additional remainder control, and an estimated natural reference law contributes an additional sampling component.

\begin{proposition}[Optimal mediator randomization at a fixed history]
Fix $(x,w,l)$ and put $a_m=g(m\mid x)\sigma_m(x,l)\ge0$. Among probabilities with $r_m>0$, the infimum of $\sum_ma_m^2/r_m$ is $(\sum_ma_m)^2$, attained by $r_m=a_m/\sum_ja_j$ when every $a_m>0$. If $r_m\ge\ell>0$ and $K\ell<1$, and some $a_m>0$, the unique optimum on positive-$a_m$ coordinates and the minimum-floor allocation on zero coordinates are
\[
 r_m=\max\{\ell,a_m/\lambda\},\qquad
 \sum_m\max\{\ell,a_m/\lambda\}=1.
\]
\end{proposition}
\begin{proof}
Cauchy--Schwarz gives $(\sum_ma_m)^2\le(\sum_ma_m^2/r_m)(\sum_mr_m)$. Equality requires $r_m$ proportional to $a_m$. For floor constraints the objective is convex and differentiable on the feasible simplex. On any coordinate above its floor, the multiplier equation is $-a_m^2/r_m^2+\lambda^2=0$. At a floor the derivative condition requires $a_m/\lambda\le\ell$. These are exactly the displayed formula and the complementary-slackness conditions, sufficient by convexity. The normalization function is continuous, decreases from infinity to $K\ell$, and is strictly decreasing whenever above $K\ell$, so its root exists and is unique. A zero-$a_m$ coordinate must be at its floor, since otherwise moving mass from it to a positive coordinate strictly improves the objective. If all $a_m$ vanish, every feasible allocation is optimal.
\end{proof}
A per-history linear cost constraint can be added. Under a strictly feasible cost bound, its multiplier $\nu\ge0$ changes the free-coordinate formula to $r_m=a_m/\sqrt{\lambda+\nu c_m}$, with floor truncation and complementary slackness; this statement presumes the denominators are positive on free coordinates. The proof is the same convex first-order inequality. It is an oracle allocation rule, not a theorem that its unknown variances can be learned costlessly.

\begin{proposition}[Randomizing the mediator does not identify a natural effect]
There are two models with identical natural-arm data and identical distributions under every sequential mediator randomization, with smooth controlled means and positive natural mediator probabilities, but with
$E[Y(1,M(0))]$ equal to $3/4$ and $1/4$, respectively.
\end{proposition}
\begin{proof}
Let $U,V_0,V_1,R$ be independent uniform variables. Set $L(0)=V_0$ and $M(0)=1\{U>1/2\}$. In model A set $L(1)=U$; in model B set $L(1)=1-U$. In both models let the natural $M(1)=1\{V_1>1/2\}$. Define
\[
 Y(0,m)=1\{R\le1/2\},\qquad
 Y(1,m)=1\{R\le mL(1)+(1-m)(1-L(1))\}.
\]
In arm zero, $L$ is uniform and independent of the fair mediator, and the outcome is a fair coin. In arm one, $L$ is uniform, the natural mediator is an independent fair coin, and the conditional outcome mean is the common linear function $ml+(1-m)(1-l)$. All observed natural-arm laws therefore coincide. Under any randomized mediator rule depending on observed $(X,W,L)$ and external randomization, the same conditional outcome law applies, so the experimental laws coincide as well.

In model A, substituting $M(0)$ into the treated response yields conditional mean $\max(U,1-U)$, whose integral is $3/4$. In model B it yields $\min(U,1-U)$, whose integral is $1/4$. In contrast the stochastic intervention drawing a fresh fair mediator gives mean $1/2$ in both models. The intermediate-covariate densities are uniform, all controlled means are smooth, and mediator positivity holds; the difference is exclusively in the unobserved cross-world dependence.
\end{proof}

\section*{Remaining gap}
The allocation formula minimizes one variance component with supplied reference law and response variances. Joint learning of reference distributions, variance functions, recruitment, costs, and treatment probabilities is unresolved here. The counterexample proves that sequential randomization cannot alone identify the natural cross-world target; an additional restriction connecting $M(w')$ to the treated response object is necessary.

\section{Completion Estimate}
45\%

This is a post hoc, heuristic estimate for the full catalogue target.
The attempt establishes stochastic mediation identification, an orthogonal variance decomposition and oracle mediator allocation, while proving natural-effect nonidentification. Joint cost-constrained treatment and mediator design with learned reference laws and variance functions, adaptive allocation and uniform efficient inference remains unresolved.

\begin{thebibliography}{9}
\bibitem{agenda} C. Cinelli, A. Feller, G. Imbens, E. Kennedy, S. Magliacane, and J. Zubizarreta. \emph{Challenges in Statistics: A Dozen Challenges in Causality and Causal Inference} (2025). \url{https://arxiv.org/html/2508.17099v1}.
\end{thebibliography}
\end{document}
